\documentclass[10pt,a4paper]{amsart}
\usepackage[margin=19mm]{geometry}
\usepackage{amsmath,amssymb,mathtools}
\usepackage[scr=rsfs,cal=euler]{mathalfa}
\usepackage{xfrac}
\usepackage[unicode,hidelinks,bookmarks=false]{hyperref}
\newcommand{\ie}{i.e.\ }
\newcommand{\cf}{cf.\ }
\newcommand{\resp}{respectively\ }
\newcommand{\Subsection}{\subsection}
\providecommand{\nobreakdash}{\nobreak\hbox{-}}
\setlength{\emergencystretch}{2em}
\multlinegap0pt
\usepackage{enumerate}
\usepackage{amssymb}
\usepackage{tikz}
\usepackage{mathtools}
\newtheorem{theorem}{Theorem}
\title{This is the title}
\author{K. Mahesh Krishna}
\date{Source: arXiv:2601.04228v1 (CC BY 4.0)}
\begin{document}
\maketitle

\noindent\emph{Source and licence.} This is the title. Version arXiv:2601.04228v1 (CC BY 4.0), distributed by arXiv under the Creative Commons Attribution 4.0 International licence, http://creativecommons.org/licenses/by/4.0/, as recorded in the arXiv licence metadata for this exact version and shown on the abstract page https://arxiv.org/abs/2601.04228v1. Full licence notice: \texttt{licenses/arxiv-2601.04228-CC-BY-4.0.txt}.\par\smallskip

\noindent\textit{Editorial scope.} Counted unit: the article's theorem NAB (source0.tex lines 554--564). The article's complete original proof of it is delivered with it (source0.tex lines 565--629, one \texttt{proof} environment of 65 lines). No mathematics is written, repaired or re-derived: the statement and the proof are the article's own lines, in the article's own notation, and no retained line is substituted or re-flowed.  No further source-local context is needed: every cross-reference of every retained line resolves inside the two delivered spans.  Nothing was omitted from any retained span, so the package carries no omission record.  No retained line contains a citation, so the wrapper carries no bibliography and no bibliography entry.  No retained line contains a figure environment, an image-inclusion command or drawing code, so no image is delivered and the package declares no asset dependency.  Editorial formatting only, in addition: an AMS \texttt{amsart} preamble that restates the article's own declarations and macros used by the retained lines, namely the environments \texttt{theorem} on separate counters, together with the standard packages those lines need.  The retained environments print in this document as Theorem 1 in order of appearance, whereas the article prints the same items with the numbering of its own full document.  The complete native source of the article as distributed in the arXiv e-print for this exact version is bundled unchanged as \texttt{sources/192281/source0.tex}.
\begin{theorem}  \label{NAB} (\textbf{Non-Archimedean Brauer Oval (of Cassini) Theorem})

For every $A=[a_{j,k}]_{1\leq j\leq n, 1\leq k \leq n}\in  \mathbb{M}_n(\mathbb{K})$,
\begin{align*}
	\sigma(A)\subseteq \bigcup_{j,k=1, j \neq k}^n\{z \in \mathbb{K}: |z-a_{j, j}||z-a_{k,k}|\leq h_j(A)h_k(A)\}
\end{align*}
and 
\begin{align*}
	\sigma(A)\subseteq \bigcup_{j,k=1, j \neq k}^n\{z \in \mathbb{K}: |z-a_{j, j}||z-a_{k,k}|\leq v_j(A)v_k(A)\}.
\end{align*}	
\end{theorem}

\begin{proof}
Let $\lambda \in \sigma(A)$. Then there exists a $0\neq \textbf{x} =(x_j)_{j=1}^n \in \mathbb{K}^n$ such that 
\begin{align}\label{PE1}
	\lambda \textbf{x} =A\textbf{x}.
\end{align}
Choose $1\leq j \leq n$ such that 
\begin{align*}
	 |x_j|=\max_{1\leq l \leq n}|x_l|.
\end{align*}
Now choose $1\leq k \leq n$ with $k \neq j$ such that 
\begin{align*}
	|x_k|=\max_{1\leq l \leq n, l\neq j}|x_l|.
\end{align*}
Then we have $1\leq j, k \leq n$ with $j \neq k$ and 
\begin{align}\label{PE2}
	|x_j|\geq |x_k|\geq \max_{1\leq l \leq n, l\neq j, l \neq k}|x_l|.
\end{align}
We have two cases. Case (i): $ |x_k|=0$. Considering the $j$-th coordinate in Equation (\ref{PE1}) gives 
\begin{align*}
	\lambda x_j=\sum_{p=1}^na_{j,p}x_p.
\end{align*}
Rewriting previous equation gives 
\begin{align*}
	(\lambda -a_{j,j})x_j=\sum_{p=1, p\neq j}^na_{j,p}x_p.
\end{align*}
Therefore using (\ref{PE2}) we get
\begin{align*}
	|(\lambda -a_{j,j})x_j|&=\left|\sum_{p=1, p\neq j}^na_{j,p}x_p\right|\leq \max_{1\leq p \leq n, p \neq j}|a_{j,p}x_p|\\
	&\leq \left(\max_{1\leq p \leq n, p \neq j}|a_{j,p}|\right)\left(\max_{1\leq p \leq n, p \neq j}|x_p|\right)=\left(\max_{1\leq p \leq n, p \neq j}|a_{j,p}|\right)|x_k|\\
	&=0.
\end{align*}
Since $	|x_j|\geq |x_k|$ for all $1\leq k \leq n$,  $ |x_k|=0$ and $\textbf{x}\neq 0$, we must have $|x_j|\neq 0$. Previous inequality then gives $	|\lambda -a_{j,j}|=0$. So 
\begin{align*}
\lambda =a_{j,j} \in \bigcup_{p, q=1, p \neq q}^n\{z \in \mathbb{K}: |z-a_{p,p}||z-a_{q,q}|\leq h_p(A)h_q(A)\}.
\end{align*}
Case (ii): $ |x_k|>0$. Considering $j$-th and $k$-th coordinates in Equation (\ref{PE1}) give
\begin{align}\label{PE3}
	(\lambda -a_{j,j})x_j=\sum_{p=1, p\neq j}^na_{j,p}x_p
\end{align}
and 
\begin{align}\label{PE4}
	(\lambda -a_{k,k})x_k=\sum_{q=1, q\neq k}^na_{k,q}x_q.
\end{align}
Multiplying Equations (\ref{PE3}) and (\ref{PE4}) and taking non-Archimedean valuation gives
\begin{align*}
|(\lambda -a_{j,j})x_j	(\lambda -a_{k,k})x_k|&=\left|\sum_{p=1, p\neq j}^na_{j,p}x_p\right|\left|\sum_{q=1, q\neq k}^na_{k,q}x_q\right|\\
&\leq \left(\max_{1\leq p \leq n, p \neq j}|a_{j,p}x_p|\right)\left(\max_{1\leq q \leq n, q \neq k}|a_{k,p}x_q|\right)\\
&\leq \left(\max_{1\leq p \leq n, p \neq j}|a_{j,p}|\right)\left(\max_{1\leq p \leq n, p \neq j}|x_p|\right)\left(\max_{1\leq q \leq n, q \neq k}|a_{k,p}|\right)\left(\max_{1\leq q \leq n, q \neq k}|x_q|\right)\\
&=\left(\max_{1\leq p \leq n, p \neq j}|a_{j,p}|\right)|x_k|\left(\max_{1\leq q \leq n, q \neq k}|a_{k,p}|\right)|x_j|\\
&=h_j(A)|x_k|h_k(A)|x_j|.
\end{align*}
Therefore 
\begin{align*}
|(\lambda -a_{j,j})(\lambda -a_{k,k})||x_jx_k|\leq h_j(A)h_k(A)|x_j||x_k|.
\end{align*}
Since $|x_jx_k|\neq0$, we have 
\begin{align*}
|(\lambda -a_{j,j})(\lambda -a_{k,k})|\leq h_j(A)h_k(A).	
\end{align*}
Previous inequality says that 
\begin{align*}
	\lambda  \in \bigcup_{p, q=1, p \neq q}^n\{z \in \mathbb{K}: |z-a_{p,p}||z-a_{q,q}|\leq h_p(A)h_q(A)\}.
\end{align*}
Second inclusion in the statement follows by considering the transpose of $A$ and noting that the spectrum of a matrix and its transpose are equal.	
\end{proof}

\end{document}
